\subsection{Lax-idempotence}
\label{subsec:2dloc-kz}

\begin{defi}
\label{def:lax-idempotent}
Let $T$ be an enriched monad on an $(\infty,2)$-category $\D$ with unit $\eta$ and multiplication $\mu$, as in (I4).
$T$ is \emph{lax-idempotent} if for every $T$-algebra $(d,a\colon Td \to d)$ the structure map is left adjoint to the unit, $a \dashv \eta_{d}$, with counit the identification $a\eta_{d} \simeq \mathrm{id}_{d}$.
It is \emph{colax-idempotent} if instead $\eta_{d} \dashv a$ with unit $\mathrm{id}_{d} \simeq a\eta_{d}$.
\end{defi}

By \cref{def:adjunction} this is a condition in $\htwo\D$.
Lax-idempotent monads are Kock's KZ-doctrines \cite{kock1995monads}, with the terminology of Kelly--Lack \cite{kellylack1997propertylike}; with our convention the free cocompletion monad on $\Cat$ is lax-idempotent, since the structure map $\colim\colon \Free(\C) \to \C$ of a cocomplete $\C$ is left adjoint to the Yoneda embedding. Idempotent monads are both lax- and colax-idempotent.

Throughout this subsection $\D \in \tprl$, $S$, $\Pc = \Pc_{S}$, $L \dashv \iota$ and $T = \iota L$ are as in \cref{thm:2dloc-local}. For $p \in \Pc$ we write $a_{p} := \iota(\varepsilon_{p})\colon T\iota p \to \iota p$, and for $h\colon d \to \iota q$ we write $h^{\sharp} := a_{q}T(h)$ for the $S$-local $1$-morphism $Td \to \iota q$ corresponding to $h$ under $\Hom_{\Pc}(Ld,q) \simeq \Hom_{\D}(d,\iota q)$; thus $h^{\sharp}\eta_{d} \simeq h$.

\begin{lem}
\label{lem:kz-criteria}
The following are equivalent.
\begin{enumerate}
\item[(a)] $T$ is lax-idempotent.
\item[(b)] For every $p \in \Pc$ there is a $2$-morphism $\theta_{p}\colon \mathrm{id}_{T\iota p} \Rightarrow \eta_{\iota p}a_{p}$ with $\theta_{p}\eta_{\iota p} = \mathrm{id}_{\eta_{\iota p}}$ and $a_{p}\theta_{p} = \mathrm{id}_{a_{p}}$.
\item[(c)] For every $d \in \D$ there is a $2$-morphism $\delta_{d}\colon T(\eta_{d}) \Rightarrow \eta_{Td}$, natural in $d$, with $\delta_{d}\eta_{d} = \mathrm{id}$ and $\mu_{d}\delta_{d} = \mathrm{id}$.
\end{enumerate}
\end{lem}

\begin{proof}
All three are statements about the pseudomonad induced by $T$ on the bicategory $\htwo\D$.
(a)$\Leftrightarrow$(b): by \cref{thm:2dloc-local}(2) every $T$-algebra is of the form $(\iota p,a_{p})$, and an adjunction $a_{p} \dashv \eta_{\iota p}$ with counit the identity is the same as a unit $\theta_{p}$ satisfying the two triangle identities, which are the two equations in (b).
(a)$\Leftrightarrow$(c) is Kock's characterization of KZ-doctrines \cite{kock1995monads}, in the form for pseudomonads of \cite{marmolejo1997doctrines}: (c) is equivalent to $T\eta_{d} \dashv \mu_{d}$, and to $\mu_{d} \dashv \eta_{Td}$, i.e.\ to (a) for the free algebras, which implies (a) for all algebras.
\todo{Check the orientation conventions in \cite{kock1995monads,marmolejo1997doctrines} against ours.}
\end{proof}

Note that in (b) and (c) the source of the $2$-morphism is $S$-local ($\mathrm{id}_{T\iota p}$, resp.\ $T(\eta_{d}) = (\eta_{Td}\eta_{d})^{\sharp}$), while its target ($\eta_{\iota p}a_{p}$, resp.\ $\eta_{Td}$) is not. The universal property of $L$ only controls $2$-morphisms between $S$-local $1$-morphisms, so neither (b) nor (c) is a formal consequence of \cref{thm:2dloc-local}. Indeed, the $\infty$-categories with finite products also form a locally full sub-$(\infty,2)$-category of $\Cat$, monadic with property-like algebras, but the corresponding monad is colax-idempotent. What distinguishes our situation is that the new $1$-morphisms of $\D\radj{S}$ are \emph{right adjoints} of the $s \in S$.

\begin{rmk}
\label{rmk:kz-coreflection}
If $T$ is lax-idempotent, then for all $d \in \D$ and $q \in \Pc$ the full subcategory $\Hom_{\Pc}(Ld,q) \subset \Hom_{\D}(Td,\iota q)$ is coreflective, with coreflection $g \mapsto (g\eta_{d})^{\sharp}$ and counit $c_{g}\colon (g\eta_{d})^{\sharp} = a_{q}T(g)T(\eta_{d}) \xRightarrow{a_{q}T(g)\delta_{d}} a_{q}T(g)\eta_{Td} \simeq a_{q}\eta_{\iota q}g \simeq g$.
Thus every $1$-morphism $g$ out of a free object has a best $S$-local approximation, the $S$-local extension of its restriction to $d$. The identities $c_{g}\eta_{d} = \mathrm{id}$ and $c_{f} = \mathrm{id}$ for $S$-local $f$ follow from (c); since $c$ is a natural transformation of functors $\Hom_{\D}(Td,\iota q) \to \Hom_{\D}(Td,\iota q)$, the criterion \cite[5.2.2.8]{htt} upgrades them to the statement on mapping spaces.
\end{rmk}

\begin{conj}
\label{conj:2dloc-kz}
The monad $T = \iota L$ of \cref{thm:2dloc-local} is lax-idempotent. Dually, the monad obtained by freely adjoining \emph{left} adjoints to $S$ is colax-idempotent.
\end{conj}

The second statement follows from the first by \cref{rmk:adjoin-duality}. We collect the evidence.

\begin{exmp}
\label{ex:kz-initial}
Let $\D = \Cat$ and $S = \{\emptyset \to \mathrm{pt}\}$. Then $s^{*}\colon \func(\mathrm{pt},\C) \to \func(\emptyset,\C) = \mathrm{pt}$ has a left adjoint if and only if $\C$ has an initial object, and a functor is $S$-local if and only if it preserves initial objects. The reflection is $L\C = \C^{\triangleleft}$ with $\eta$ the inclusion, and for $\C \in \Pc$ the structure map $a\colon \C^{\triangleleft} \to \C$ sends the cone point $-\infty$ to the initial object $\emptyset$. Then $a \dashv \eta$: $\Hom_{\C}(a(-\infty),c) = \Hom_{\C}(\emptyset,c) \simeq \mathrm{pt} \simeq \Hom_{\C^{\triangleleft}}(-\infty,c)$, and $\theta$ is the unique morphism $-\infty \to \emptyset$. So \cref{conj:2dloc-kz} holds here.
\end{exmp}

\begin{rmk}[the universal case]
\label{rmk:kz-universal}
For $\D = \Free(\A)$, $L = j_{!}$ and $\iota = j^{*}$ for the lax epimorphism $j\colon \A \to \B := \A\radj{S}$, and $T = j^{*}j_{!}$ preserves colimits. So $\delta$ in \cref{lem:kz-criteria}(c) is determined by its values on representables, where it is a family of morphisms $[m,\mathrm{id}_{ja}] \to [\mathrm{id}_{ja'},m]$ in the coends $\int^{a''}\Hom_{\B}(ja',ja'') \times \Hom_{\B}(ja'',ja)$, natural in $m \in \Hom_{\B}(ja',ja)$.
For $\A = \Delta^{1}$, $\B = \Adj$ and $m = r$, such a morphism is
\[
[r,\mathrm{id}_{-}] \xrightarrow{(r,\eta)} [r,rl] \simeq [lr,r] \xrightarrow{(\varepsilon,r)} [\mathrm{id}_{+},r],
\]
using the coend relation for $l = j(0 \to 1)$. It uses the unit and the counit in the only direction in which they exist, and this is where the lax, rather than colax, direction comes from. Turning this into a proof requires a description of the hom-categories of $\A\radj{S}$, e.g.\ by the zig-zag formula of \cite{abellan2026bifibrations}, and then passing from presheaf categories to general $\D \in \tprl$.
\todo{Prove \cref{conj:2dloc-kz}.}
\end{rmk}

\begin{rmk}
\label{rmk:kz-not-pointwise}
For $\D = \Cat$ one might try to identify $L$ with a free cocompletion under the colimits computing left Kan extensions along the $s \in S$, and deduce \cref{conj:2dloc-kz} from the lax-idempotence of free cocompletions. This fails: $S$-local objects are those admitting (global) left Kan extensions, which need not be pointwise, see \cref{ex:adjoin-basic}(3).
\end{rmk}

\begin{rmk}
\label{rmk:2dloc-two-sided}
If one adjoins right adjoints to a set $S_{1}$ and left adjoints to a set $S_{2}$ simultaneously (pushout along $\Free(\ell)$ on the $S_{1}$-summands and $\Free(\rho)$ on the $S_{2}$-summands), \cref{thm:2dloc-local} goes through verbatim, with both Beck--Chevalley conditions of \cref{lem:mates}, and the resulting monad is property-like in the sense of Kelly--Lack. It is in general neither lax- nor colax-idempotent, just as the monad ``free initial and terminal object'' on $\Cat$ is not.
\end{rmk}

\begin{rmk}
\label{rmk:2dloc-6ff}
For an $(\infty,1)$-category $\C$ with classes $I$, $P$ of morphisms, the universal property of six-functor formalisms of Cnossen--Lenz--Linskens \cite{cnossenlenzlinskens2025universal} says that a six-functor formalism on $(\C,I,P)$ is a functor out of $(\C\opcat)\radj{I}\ladj{P}$ (in $\prl$-valued form) in which, in addition, the Beck--Chevalley $2$-morphisms for the pullback squares of $\C$ are invertible. The first step is the two-dimensional localization of this section; the second is an ordinary (idempotent) localization at a set of $2$-morphisms. This is taken up in the next section.
\todo{Write the next section.}
\end{rmk}
